%%%%%%%%%%%%%%%%%%%%%%%%%%%%%%%%%%%%%%%%%%%%%%%%%%%%%%%%%%%%%%%%%%%%%%
% Lớp 9 | Chương V. Đường tròn | Vị trí tương đối của đường thẳng và đường tròn
% Dạng 9H5?16-2: Tiếp tuyến: dấu hiệu nhận biết, chứng minh tiếp tuyến
%%%%%%%%%%%%%%%%%%%%%%%%%%%%%%%%%%%%%%%%%%%%%%%%%%%%%%%%%%%%%%%%%%%%%%


% nguồn: 2026-11-21_lop9_chuong5_bai14_17 (nhập 21/11/2026)
\begin{ex}%[9H5N16-2]%[ma:8ade643d]
% nguon: SGK Toán 9 KNTT tập 1 – Bài 16, Vị trí tương đối của đường thẳng và đường tròn
	Đường thẳng $d$ là tiếp tuyến của đường tròn $(O; R)$ tại điểm $A$ khi và chỉ khi
	\choice
	{$A$ nằm trên đường tròn $(O)$}
	{$d$ đi qua tâm $O$}
	{\True $A$ nằm trên đường tròn $(O)$ và $d \perp OA$ tại $A$}
	{$OA = 2R$ và $d \perp OA$ tại $A$}
	\loigiai{
		Theo dấu hiệu nhận biết tiếp tuyến, đường thẳng $d$ là tiếp tuyến của đường tròn $(O; R)$ tại điểm $A$ khi và chỉ khi $A \in (O)$ và $d \perp OA$ tại $A$.
	}
\end{ex}

% nguồn: 2026-11-21_lop9_chuong5_bai14_17 (nhập 21/11/2026)
\begin{ex}%[9H5H16-2]%[ma:a7b990e8]
% nguon: SGK Toán 9 KNTT tập 1 – Bài 16, Vị trí tương đối của đường thẳng và đường tròn
	Cho tam giác $ABC$ có $AB = 6\text{ cm}$, $AC = 8\text{ cm}$, $BC = 10\text{ cm}$. Xét đường tròn tâm $C$ bán kính $R = 8\text{ cm}$. Xét tính đúng, sai của các khẳng định sau:
	\choiceTF
	{\True Tam giác $ABC$ vuông tại $A$}
	{\True Điểm $A$ nằm trên đường tròn $(C; 8\text{ cm})$}
	{Đường thẳng $AB$ cắt đường tròn $(C; 8\text{ cm})$ tại hai điểm phân biệt}
	{\True Khoảng cách từ điểm $A$ đến đường thẳng $BC$ bằng $4{,}8\text{ cm}$}
	\loigiai{
		\begin{itemchoice}
			\itemch Đúng. Ta có $AB^2 + AC^2 = 6^2 + 8^2 = 100 = 10^2 = BC^2$, theo định lí Pythagore đảo thì tam giác $ABC$ vuông tại $A$.
			\itemch Đúng. Vì $CA = 8\text{ cm} = R$ nên điểm $A$ thuộc đường tròn $(C; 8\text{ cm})$.
			\itemch Sai. Vì $A \in (C; 8\text{ cm})$ và $AB \perp CA$ tại $A$ nên đường thẳng $AB$ là tiếp tuyến của đường tròn $(C; 8\text{ cm})$, chỉ có một điểm chung là $A$.
			\itemch Đúng. Đường cao $AH$ của tam giác vuông $ABC$ thoả mãn $AH \cdot BC = AB \cdot AC \Rightarrow AH = \dfrac{6 \cdot 8}{10} = 4{,}8\text{ cm}$.
		\end{itemchoice}
	}
\end{ex}

% nguồn: 2026-11-21_lop9_chuong5_bai14_17 (nhập 21/11/2026)
\begin{bt}%[9H5V16-2]%[ma:fdaf28b6]
% nguon: SGK Toán 9 KNTT tập 1 – Bài 16, Vị trí tương đối của đường thẳng và đường tròn
	Cho đường tròn $(O; 5\text{ cm})$. Từ điểm $M$ nằm ngoài đường tròn sao cho $OM = 13\text{ cm}$, kẻ tiếp tuyến $MA$ đến đường tròn $(O)$ với $A$ là tiếp điểm. Tính độ dài đoạn thẳng $MA$ (đơn vị cm).
	\shortans{12}
	\loigiai{
		Vì $MA$ là tiếp tuyến của $(O)$ tại $A$ nên $MA \perp OA$, suy ra tam giác $OAM$ vuông tại $A$.\\
		Áp dụng định lí Pythagore: $MA = \sqrt{OM^2 - OA^2} = \sqrt{13^2 - 5^2} = \sqrt{144} = 12\text{ cm}$.
	}
\end{bt}

% nguồn: 2026-11-21_lop9_chuong5_bai16_17_bs (nhập 21/11/2026)
\begin{ex}%[9H5N16-2]%[ma:fc30a7f8]
% nguon: SGK Toán 9 KNTT tập 1 – Bài 16, Vị trí tương đối của đường thẳng và đường tròn
	Nếu đường thẳng $a$ là tiếp tuyến của đường tròn $(O)$ tại điểm $M$ thì khẳng định nào sau đây đúng?
	\choice
	{$a$ đi qua tâm $O$}
	{\True $a \perp OM$ tại $M$}
	{$a$ song song với $OM$}
	{$OM$ là đường kính của $(O)$}
	\loigiai{
		Theo tính chất tiếp tuyến của đường tròn, nếu đường thẳng $a$ tiếp xúc với đường tròn $(O)$ tại $M$ thì $a$ vuông góc với bán kính $OM$ tại tiếp điểm $M$.
	}
\end{ex}

% nguồn: 2026-11-21_lop9_chuong5_bai16_17_bs (nhập 21/11/2026)
\begin{ex}%[9H5N16-2]%[ma:df79c9e0]
% nguon: SGK Toán 9 KNTT tập 1 – Bài 16, Vị trí tương đối của đường thẳng và đường tròn
	Qua một điểm $A$ nằm trên đường tròn $(O)$, ta có thể kẻ được bao nhiêu tiếp tuyến với đường tròn $(O)$?
	\choice
	{$0$}
	{\True $1$}
	{$2$}
	{Vô số}
	\loigiai{
		Qua một điểm $A$ nằm trên đường tròn $(O)$ chỉ có duy nhất một đường thẳng vuông góc với bán kính $OA$ tại $A$, nên chỉ kẻ được đúng $1$ tiếp tuyến.
	}
\end{ex}

% nguồn: 2026-11-21_lop9_chuong5_bai16_17_bs (nhập 21/11/2026)
\begin{ex}%[9H5H16-2]%[ma:5f800d15]
% nguon: SGK Toán 9 KNTT tập 1 – Bài 16, Vị trí tương đối của đường thẳng và đường tròn
	Cho đường tròn $(O; 6\text{ cm})$ và điểm $M$ nằm ngoài đường tròn sao cho $OM = 10\text{ cm}$. Kẻ tiếp tuyến $MA$ đến $(O)$ với $A$ là tiếp điểm. Xét tính đúng, sai của các khẳng định sau:
	\choiceTF
	{\True Tam giác $OAM$ vuông tại $A$}
	{\True Độ dài đoạn thẳng $MA$ bằng $8\text{ cm}$}
	{Diện tích tam giác $OAM$ bằng $48\text{ cm}^2$}
	{\True Đường cao $AH$ kẻ từ $A$ xuống $OM$ có độ dài bằng $4{,}8\text{ cm}$}
	\loigiai{
		\begin{itemchoice}
			\itemch Đúng. Vì $MA$ là tiếp tuyến của $(O)$ tại $A$ nên $MA \perp OA$, do đó tam giác $OAM$ vuông tại $A$.
			\itemch Đúng. Áp dụng định lí Pythagore: $MA = \sqrt{OM^2 - OA^2} = \sqrt{10^2 - 6^2} = 8\text{ cm}$.
			\itemch Sai. Diện tích tam giác vuông $OAM$ là $\dfrac{1}{2} \cdot OA \cdot MA = \dfrac{1}{2} \cdot 6 \cdot 8 = 24\text{ cm}^2$.
			\itemch Đúng. Ta có $AH \cdot OM = OA \cdot MA \Rightarrow AH = \dfrac{6 \cdot 8}{10} = 4{,}8\text{ cm}$.
		\end{itemchoice}
	}
\end{ex}

% nguồn: 2026-11-21_lop9_chuong5_bai16_17_bs (nhập 21/11/2026)
\begin{ex}%[9H5H16-2]%[ma:15fb1f99]
% nguon: SGK Toán 9 KNTT tập 1 – Bài 16, Vị trí tương đối của đường thẳng và đường tròn
	Cho tam giác $ABC$ vuông tại $A$ có $AB = 5\text{ cm}$, $AC = 12\text{ cm}$. Xét đường tròn tâm $B$ bán kính $R = 5\text{ cm}$. Xét tính đúng, sai của các khẳng định sau:
	\choiceTF
	{\True Điểm $A$ nằm trên đường tròn $(B; 5\text{ cm})$}
	{\True Đường thẳng $AC$ là tiếp tuyến của đường tròn $(B; 5\text{ cm})$}
	{Độ dài cạnh $BC$ bằng $15\text{ cm}$}
	{\True Điểm $C$ nằm ngoài đường tròn $(B; 5\text{ cm})$ và khoảng cách từ $C$ đến điểm gần nhất trên đường tròn $(B; 5\text{ cm})$ bằng $8\text{ cm}$}
	\loigiai{
		\begin{itemchoice}
			\itemch Đúng. Vì $BA = 5\text{ cm} = R$ nên $A$ thuộc đường tròn $(B; 5\text{ cm})$.
			\itemch Đúng. Vì $A \in (B; 5\text{ cm})$ và $AC \perp BA$ tại $A$ nên $AC$ là tiếp tuyến của $(B; 5\text{ cm})$.
			\itemch Sai. Trong tam giác $ABC$ vuông tại $A$: $BC = \sqrt{5^2 + 12^2} = 13\text{ cm}$.
			\itemch Đúng. Vì $BC = 13 > 5$ nên $C$ nằm ngoài $(B; 5\text{ cm})$, khoảng cách ngắn nhất từ $C$ đến đường tròn là $BC - R = 13 - 5 = 8\text{ cm}$.
		\end{itemchoice}
	}
\end{ex}

% nguồn: 2026-11-21_lop9_chuong5_bai16_17_bs (nhập 21/11/2026)
\begin{bt}%[9H5V16-2]%[ma:e3496959]
% nguon: SGK Toán 9 KNTT tập 1 – Bài 16, Vị trí tương đối của đường thẳng và đường tròn
	Cho đường tròn $(O; 8\text{ cm})$. Từ điểm $P$ nằm ngoài đường tròn kẻ tiếp tuyến $PA$ ($A$ là tiếp điểm) sao cho $PA = 15\text{ cm}$. Tính độ dài đoạn thẳng $OP$ (đơn vị cm).
	\shortans{17}
	\loigiai{
		Vì $PA$ là tiếp tuyến của $(O)$ tại $A$ nên tam giác $OAP$ vuông tại $A$.\\
		Theo định lí Pythagore: $OP = \sqrt{OA^2 + PA^2} = \sqrt{8^2 + 15^2} = \sqrt{289} = 17\text{ cm}$.
	}
\end{bt}

% nguồn: 2026-11-21_lop9_chuong5_bai16_17_bs (nhập 21/11/2026)
\begin{bt}%[9H5V16-2]%[ma:97b3a3c2]
% nguon: SGK Toán 9 KNTT tập 1 – Bài 16, Vị trí tương đối của đường thẳng và đường tròn
	Từ điểm $M$ nằm ngoài đường tròn $(O; 9\text{ cm})$ sao cho $OM = 15\text{ cm}$, kẻ tiếp tuyến $MA$ đến đường tròn $(O)$ với $A$ là tiếp điểm. Tính diện tích tam giác $OAM$ (đơn vị $\text{cm}^2$).
	\shortans{54}
	\loigiai{
		Tam giác $OAM$ vuông tại $A$ nên $MA = \sqrt{OM^2 - OA^2} = \sqrt{15^2 - 9^2} = 12\text{ cm}$.\\
		Diện tích tam giác $OAM$ là $S = \dfrac{1}{2} \cdot OA \cdot MA = \dfrac{1}{2} \cdot 9 \cdot 12 = 54\text{ cm}^2$.
	}
\end{bt}
